\documentclass[12pt]{article}
\usepackage[margin=0.75in]{geometry}
\usepackage{fontspec}
\setmainfont{Latin Modern Roman}
\usepackage{microtype}
\usepackage{multicol}
\usepackage{fancyhdr}
\usepackage{titlesec}
\usepackage{enumitem}
\usepackage{xcolor}
\usepackage[hidelinks]{hyperref}
\setlength{\columnsep}{0.28in}
\setlength{\parindent}{1.1em}
\setlength{\parskip}{0.32em}
\setlength{\headheight}{15pt}
\titleformat{\section}{\large\bfseries\scshape}{}{0pt}{}
\titleformat{\subsection}{\normalsize\bfseries}{}{0pt}{}
\pagestyle{fancy}
\fancyhf{}
\fancyfoot[C]{\thepage}
\renewcommand{\headrulewidth}{0.4pt}
\newcommand{\focusbox}[1]{\par\vspace{0.4em}\noindent\begingroup\setlength{\fboxsep}{6pt}\colorbox{gray!12}{\begin{minipage}{\dimexpr\linewidth-2\fboxsep\relax}#1\end{minipage}}\endgroup\par\vspace{0.5em}}

\fancyhead[L]{Whately: Circular Reasoning}
\fancyhead[R]{Student Reading}
\begin{document}
\begin{center}
{\LARGE\bfseries Circular Reasoning and Begging the Question}\par\vspace{0.25em}
{\large\scshape Week 9: When the Proof Secretly Assumes the Answer}\par\vspace{0.5em}
{Adapted and modernized from Richard Whately, \textit{Elements of Logic}}\par\vspace{0.6em}
\rule{0.82\textwidth}{0.4pt}
\end{center}

\section*{Central Question}
How can an argument sound like proof while merely repeating its conclusion?

\section*{Vocabulary Preview}
\begin{multicols}{2}
\begin{itemize}[leftmargin=*, itemsep=0.18em]
\item \textbf{Circular reasoning}: an argument whose reason depends on already accepting the conclusion.
\item \textbf{Begging the question}: assuming in the premises the very point that needs proof (petitio principii).
\item \textbf{Assumed conclusion}: the claim that is treated as settled before independent support is given.
\item \textbf{Independent support}: a reason that could be accepted without already granting the conclusion.
\item \textbf{Hidden premise}: an unstated assumption that the argument needs.
\item \textbf{Restatement}: saying the conclusion again in different words and calling it a reason.
\item \textbf{Loaded definition}: a definition that smuggles the conclusion into the meaning of a key term.
\item \textbf{Question-begging epithet}: a descriptive label that assumes guilt, virtue, or proof.
\item \textbf{Rewrite}: rebuilding the argument so the reason can be tested apart from the conclusion.
\item \textbf{Petitio principii}: the traditional name for begging the question.
\item \textbf{Proof burden}: what the speaker still owes the hearer.
\item \textbf{Essay trap:} writing that restates the thesis as if restating were evidence.
\end{itemize}
\end{multicols}

\section*{Introduction}
In Week 1 you learned to find hidden assumptions. In Week 7 you learned that some arguments only \textit{look} like proof. This week names a common family of that failure: the argument that assumes what it was supposed to prove.

Whately and the older logical tradition call this \textit{petitio principii}---begging the question. The phrase does not mean ``this raises a further question'' in casual modern speech. In logic, it means the reasoning takes the conclusion for granted inside the support.

\focusbox{\textbf{Source note:} Adapted and modernized from Whately's treatment of arguments that assume what they must prove. Classroom examples are original. This week deepens Week 1's hidden-assumption work and Week 7's fallacy autopsy for a specific defect pattern.}

\begin{multicols}{2}
\section*{Reading}

\subsection*{1. What Circular Reasoning Is}
An argument is circular when the reason only works if the conclusion is already accepted. The speaker seems to move from support to claim, but the support secretly contains the claim.

Simple form:
\begin{itemize}[leftmargin=*]
\item The policy is wise because it is the smart choice.
\end{itemize}
``Smart choice'' is often just another way of saying ``wise policy.'' Nothing independent has been offered.

\subsection*{2. Restatement Is Not Evidence}
Students meet this problem constantly in essays:
\begin{itemize}[leftmargin=*]
\item Thesis: Macbeth is ambitious.
\item ``Support'': Macbeth wants power more than anything.
\end{itemize}
That may be a restatement, not a proof. A non-circular version would point to choices, speeches, or actions that a reader could accept without already granting the thesis word-for-word.

\subsection*{3. Loaded Definitions}
Sometimes the circle hides in a definition:
\begin{itemize}[leftmargin=*]
\item A true fan never criticizes the coach.
\item You criticized the coach.
\item Therefore you are not a true fan.
\end{itemize}
If ``true fan'' has been defined as ``person who never criticizes the coach,'' the conclusion was built into the definition. The argument did not discover your status; it assumed a definition that forces the result.

\subsection*{4. Question-Begging Labels}
Labels can do the same work: \textit{traitor}, \textit{real teammate}, \textit{obviously guilty}, \textit{anti-school}, \textit{natural leader}. If the label assumes the verdict, it cannot be used as neutral proof of the verdict.

\subsection*{5. How to Test for Circularity}
Ask:
\begin{enumerate}[leftmargin=*,itemsep=0.12em]
\item What is the conclusion?
\item What reason is offered?
\item Could a fair opponent accept the reason without already accepting the conclusion?
\item If not, what independent evidence would actually be needed?
\end{enumerate}
If the reason cannot be accepted first, the argument is begging the question.

\subsection*{6. Circular Versus Merely Weak}
Not every weak argument is circular. ``The cafeteria was loud; therefore the dress code is unfair'' is a bad inference, but it does not necessarily restate the conclusion. Circular reasoning is the special case where the conclusion is smuggled into the support.

Not every circular argument uses identical words. Often the circle uses a synonym, a loaded definition, or a moral title.

\subsection*{7. Why Circular Arguments Feel Strong}
They feel strong because they are coherent. They do not openly contradict themselves. They sound settled. In a speech, they can use confidence and identity pressure. But coherence is not independent proof.

\subsection*{8. Literary Cases}
In \textit{Macbeth}, manhood language can become circular: daring the murder is treated as what it means to be a man, and being a man is treated as the reason to dare the murder. In \textit{Julius Caesar}, claims about honor, conspiracy, or tyranny can assume the verdict that still needed proof. The literary task is not to accuse every character of a textbook label. The task is to ask whether a speaker's reason can be granted without already granting the conclusion.

\subsection*{9. Repair}
A repaired argument replaces the circular reason with independent support, or it narrows the claim to what has actually been shown. Example:
\begin{itemize}[leftmargin=*]
\item Circular: ``The rule is necessary because we must have this rule.''
\item Better: ``The rule is necessary because late arrivals delayed three labs last week; here is the attendance log.''
\end{itemize}

\subsection*{10. What This Week Adds}
A complete Week 9 diagnosis should include:
\begin{itemize}[leftmargin=*,itemsep=0.12em]
\item the conclusion;
\item the offered reason;
\item where the conclusion is assumed (restatement, loaded definition, or label);
\item what independent support would be needed;
\item a rewritten non-circular version or a narrowed claim.
\end{itemize}

\section*{Reading Focus}
\begin{enumerate}[leftmargin=*]
\item In logical vocabulary, what does ``begging the question'' mean?
\item How is circular reasoning different from a merely weak leap?
\item Why is restating a thesis not the same as supporting it?
\item What is a loaded definition? Give an original example.
\item List the four test questions for circularity.
\item Why can circular arguments feel coherent and still fail as proof?
\item How might manhood or honor language become circular in Shakespeare?
\end{enumerate}

\section*{Handout Summary}
Write 5--7 sentences explaining how an argument can assume its own conclusion and how a writer should repair that failure. Your summary must include the words \textit{circular}, \textit{assumption}, \textit{independent}, and \textit{rewrite}.
\end{multicols}
\end{document}
